\documentclass[sigconf,nonacm]{acmart}
\settopmatter{printacmref=false, printfolios=true}
\renewcommand\footnotetextcopyrightpermission[1]{}
\pagestyle{plain}
\usepackage{booktabs}
\usepackage{tabularx}
\usepackage{multirow}
\newcommand{\PENDING}[1]{\textbf{[#1]}}
\begin{document}
\sloppy

\title{Are LLM Written Unit Tests Flaky?}
\subtitle{Progress Report 3, ECE 382V Software Testing in the Era of Nondeterminism, Fall 2026}
\author{Shaunak Kapur (sk62288)}
\affiliation{\institution{The University of Texas at Austin}\city{Austin}\state{Texas}\country{USA}}
\author{Manas Pathak (map7848)}
\affiliation{\institution{The University of Texas at Austin}\city{Austin}\state{Texas}\country{USA}}
\maketitle

%SUMMARY%

\section{Response to the Progress Report 2 feedback}
Table~\ref{tab:response} answers each comment.

\begin{table*}[t]
\caption{Each comment on Progress Report 2 and what we changed}
\label{tab:response}
\small
\begin{tabularx}{\textwidth}{p{5.2cm}X}
\toprule
Comment & What we did \\
\midrule
The E1, E2, E3 labels were hard to follow and the parts are conceptually different. Write in numbered order. & We dropped the labels. The report now follows the four research questions of the proposal (RQ1 to RQ4) and names the two studies by what they hold fixed, the \emph{library study} and the \emph{known flaky code study}. Section~\ref{sec:setup} is a paper style first draft of the experimental setup. \\
Explain the different run modes, and make ``some runs'' concrete. & Section~\ref{sec:collect} gives each study's protocol exactly: the library study runs 300 shuffled runs, 150 with one of 10 fixed time zones, and 150 under one of 10 fixed stress-ng configurations; the known flaky code study alternates file order (even runs) and shuffled order (odd runs, seed $7000+i$) for 200 runs. It also says why the two differ. \\
The feature table looked only at order dependence. Add randomness and concurrency, and show whether human tests use these features and stay stable. & Table~\ref{tab:features} now covers randomness, time, concurrency, and shared state, and for each gives how many human and LLM test functions use it and how many of those flipped (Section~\ref{sec:rq3}). \\
Do the mitigations patch tests or improve generation? How were the prompts derived, and does that bias the evaluation? & Fixtures patch a suite after generation, prompts change generation (Section~\ref{sec:rq4}). Both were derived on the first 17 projects with GPT~4.1 mini, so we now test them on suites they were not built on: the fixtures on 656 new held out suites and 136 order dependent arm suites, and guard v3 on two new models (Section~\ref{sec:rq4}). \\
Measure quality after mitigation, for example coverage and mutation score. & New: line coverage and a fixed mutation set per project for every plain and guard suite (Table~\ref{tab:quality}). The guard prompts cost quality, mostly through suites that no longer import. \\
\bottomrule
\end{tabularx}
\end{table*}

\section{Experimental setup (first draft for the final report)}
\label{sec:setup}

\subsection{Research questions and the two studies}
We answer the four research questions of our proposal.
\begin{description}
\item[RQ1, prevalence.] How often are LLM written unit tests flaky, compared with developer written tests of the same code?
\item[RQ2, root causes.] Why do the flaky LLM tests flip?
\item[RQ3, risky features.] Do the features of a test that are known to cause flakiness (randomness, the clock, sleeps, concurrency, shared state) tell which tests flip?
\item[RQ4, mitigation.] Can a cheap step in the generate and filter pipeline remove the flaky tests, and what does it cost in test quality?
\end{description}
Two studies feed these questions. They differ in what is held fixed, so we keep them separate.
\begin{itemize}
\item \textbf{The library study} asks whether LLM tests are flaky on ordinary code. It uses 12 popular Python libraries, the developer suite of each, and eight generation conditions (Table~\ref{tab:conds}).
\item \textbf{The known flaky code study} asks whether LLM tests are flaky where flakiness is likely. It uses 41 projects in which Gruber et al.~\cite{gruber2021} found a flaky developer test, and asks four models to test the same code that test exercises.
\end{itemize}
\begin{table}[h]
\caption{Generation conditions of the library study. All generation is one shot with no test execution, except AG-S. Codes: AG agent, OS one shot, TP TestPilot style, PL plain; S Sonnet, H Haiku, G5 GPT~5.4 mini, G4 GPT~4.1 mini.}
\label{tab:conds}
\small
\begin{tabularx}{\columnwidth}{lXl}
\toprule
Code & Model & Prompt \\
\midrule
Human & library developers & none \\
AG-S & Claude Sonnet 5, agent that runs and fixes its tests & feature \\
OS-S & Claude Sonnet 5.5 & feature \\
OS-H & Claude Haiku 4.5 & feature \\
TP-S & Claude Sonnet 5.5 & TestPilot style \\
PL-S & Claude Sonnet 5.5 & plain \\
PL-G5 & GPT 5.4 mini (OpenAI) & plain \\
PL-G4 & GPT 4.1 mini (OpenAI) & plain \\
\bottomrule
\end{tabularx}
\end{table}
The \emph{feature} prompt names the time, randomness, thread, and file features of each library so the model tests them. The \emph{TestPilot style} prompt shows only signatures, docstrings, and README examples, as the TestPilot generator does~\cite{schafer2024}. The \emph{plain} prompt only asks for pytest unit tests of the given source.

Both studies use the same definitions, the same root cause method, and the same analysis. They use different run protocols because they measure different things (Section~\ref{sec:collect}).

\subsection{Subjects}
\textbf{Library study.} We took the 5,000 most downloaded PyPI packages and kept libraries (no stubs or pytest plugins) that link a GitHub repository, ship a pure Python wheel, need at most three dependencies, run on Python 3.12, and have test files (1,959 left). We split them by whether the library imports \texttt{time}, \texttt{random}, \texttt{threading}, or \texttt{asyncio}, because nondeterminism enters through those modules, and oversample that stratum on purpose. Our 12 libraries pass every rule, with 10 in the first stratum and 2 controls.

\textbf{Known flaky code study.} A project enters if its flaky developer test reproduces on our runner, which means it both passes and fails in up to 150 runs at the commit Gruber et al.\ recorded, in a fresh Python 3.8 environment. Candidates come from three fixed rules, each written before any run.
\begin{enumerate}
\item The 58 hand labeled tests of Gruber et al.\ that are not order dependent and not caused by the network (17 projects used).
\item One unlabeled test per remaining project, preferring tests not marked as infrastructure and with a failure rate near one half, and dropping names that point to the network (\texttt{scripts/select\_ext.py}, 24 projects used).
\item Two batches added after Progress Report 2, both by rules fixed before any run. The time batch takes one time, concurrency, or async wait test per project that was not tried before (\texttt{select\_time.py}, 16 eligible). The order dependent batch takes one test per project that Gruber et al.\ mark as order dependent (\texttt{select\_od.py}, first 45 of 499 eligible projects by name). An order dependent test reproduces if it both passes and fails when its whole test file runs in a new random order in each of up to 60 runs (\texttt{gruber\_reproduce.py --shuffle-file}).
\end{enumerate}
The first two rules give the 41 projects of the main known flaky code study. The order dependent batch gives a separate \emph{order dependent arm} of 17 projects that we analyze apart, because those projects flip for a different reason.
We drop a test that needs the network, a download, or special hardware.

\subsection{Test generation}
Every model writes tests one shot, with no execution, except the agentic condition AG-S. The prompts are in \texttt{prompts/README.md}. The plain prompt is a source only baseline that approximates a simple developer request, and we did not observe how developers actually prompt. In the known flaky code study every model gets the plain prompt with the target source files pasted in, and is told which code to test. The first two suites per model came from two pipelines (Claude Code subagents for the Claude models, the API for the OpenAI models). The held out sample (Section~\ref{sec:heldout}) puts every model through one pipeline, with Claude models called through \texttt{claude -p} with every tool turned off, no project settings, and the source pasted into the prompt exactly as for the OpenAI API. The order dependent arm uses the same pipeline, four suites per model. In both, each suite is a fresh call with no memory of earlier suites.

\subsection{Running tests and collecting results}
\label{sec:collect}
Every run is a fresh pytest process with a new \texttt{PYTHONHASHSEED}, and every result comes from pytest's JUnit report. A test is \emph{flaky} if it both passed and failed on the same machine. A generated test is \emph{kept} if it passes the first run of its suite, which is the one run filter of TestPilot~\cite{schafer2024}. It is \emph{removed} if it fails that run, or if it only flips after a removed test.
\begin{itemize}
\item \textbf{Library study.} The question is how often a test flips under common sources of nondeterminism, so every run shuffles the order with pytest-randomly. Each generated test runs 600 times: 300 with only the shuffle and hash seed, 150 that also set the time zone, cycling through 10 fixed zones by run index, and 150 under stress-ng load, cycling through 10 fixed CPU and memory configurations sampled once with a fixed seed, as Shaker~\cite{silva2020} does. Developer suites are about five times larger, so their tests run 200 times with the same 2:1:1 split. All runs use 2 vCPU GitHub Actions Linux runners.
\item \textbf{Known flaky code study.} The question is whether a generate and filter pipeline would ship a flaky test, and that pipeline runs the file in its written order. So run $i$ of 200 keeps the file order without pytest-randomly when $i$ is even and shuffles with seed $7000+i$ when $i$ is odd. Comparing the two halves separates order dependent tests from the rest. Because these projects are already known to flip without load, there is no time zone or load mode. The flaky developer test runs 200 times next to the generated suites as a control. We record every outcome in run order, so filters with $k$ extra runs can be measured directly.
\end{itemize}

\subsection{Root causes}
Each flaky test gets a category from Luo et al.~\cite{luo2014} and Gruber et al.~\cite{gruber2021}, based on the test, the code it calls, and a reproduction of the failure. For every flaky test in the known flaky code study we also rerun the test alone 30 times and, if it never fails alone, run every other test of its suite before it, as iDFlakies~\cite{lam2019} does, to find its polluter. Manas labels and Shaunak labels the same tests blind (without the model or the first label), and we report Cohen's kappa.

\subsection{Mitigation and test quality}
We study two kinds of mitigation. \emph{Fixtures} patch the generated suite after generation and leave the tests unchanged. M1 seeds every generator inside each test call. M2 also imports the project once at session start and restores interpreter state after each test. \emph{Prompts} change the generation itself. Guard v1 lists the usual causes of flakiness. Guard v2 adds three rules we derived from the root causes of the flaky tests that guard v1 still produced, and guard v3 narrows the seeding rule after v2 made suites fail to import. Both fixtures and the guard prompts were derived from the first 17 projects, so we evaluate them on tests and projects they were not built on (Section~\ref{sec:heldout}). For every suite we also measure quality, so a mitigation cannot win by deleting tests. We report the kept tests, the line coverage of the target files with pytest-cov, and a mutation score. The mutation score uses one fixed set of up to 40 single token mutants per project (comparison and arithmetic operators, \texttt{and}/\texttt{or}, numeric constants, booleans, and return values), sampled with a seed, so every suite of a project faces the same mutants (\texttt{scripts/e2\_quality.py}).

\subsection{Analysis}
Flaky tests cluster within suites and projects, so we compare models on suites and projects, not on tests. The main test is an exact sign test on the projects in which only one of two models has a kept flaky test, with a Bonferroni correction for six model pairs ($p < 0.0083$). For the held out sample this test was fixed in Progress Report 2 and is run unchanged from commit \texttt{7c05b6b} (\texttt{scripts/heldout\_stats.py}). Rates come with 95\% Wilson intervals.


\section{Results since Progress Report 2}
\label{sec:results}

\subsection{RQ1, prevalence: the held out sample}
\label{sec:heldout}
Progress Report 2 fixed one step before any new data: generate 8 new plain suites per model for each of the 41 projects through one pasted source pipeline, run each 200 times plain, with M1, and with M2, and analyze them alone with the frozen sign test. All 656 Claude suites ran (Table~\ref{tab:heldout}). \PENDING{GPTHELDOUT}

\begin{table}[h]
\caption{Held out sample, 41 projects, 8 new suites per model, 200 runs per suite in each column. Kept counts tests that pass the first run. Flaky kept has suites with one and projects with one in parentheses. Shuffle only counts kept flaky tests that never failed in a file order run.}
\label{tab:heldout}
\small
\setlength{\tabcolsep}{3pt}
\begin{tabularx}{\columnwidth}{Xrrrr}
\toprule
Model & Suites & Kept & Flaky kept & Shuffle only \\
\midrule
Sonnet 5.5 & 328 & 6,898 & 18 (8, 6) & 14 \\
Haiku 5.5 & 319 & 6,981 & 34 (15, 8) & 22 \\
%GPTROWS%
\bottomrule
\end{tabularx}
\end{table}

\begin{itemize}
\item \textbf{Claude models do write flaky tests once the pipeline is the same.} In Progress Report 2 Sonnet had 1 kept flaky test in 81 suites. Through the pasted source pipeline it has 18 in 328 suites (0.26\% of kept tests, 95\% Wilson interval 0.17 to 0.41), and Haiku 34 in 319 (0.49\%, 0.35 to 0.68). The share of Sonnet suites with a kept flaky test went from 1.2\% to 2.4\%, which is within noise (Fisher's exact test on suites, $p = 1.0$). So the earlier near zero for Sonnet was partly a small sample, not proof that Sonnet avoids flakiness.
\item \textbf{Most new flaky tests are order dependent.} 36 of the 52 kept flaky Claude tests never failed in a file order run, so the one run filter in file order passed them.
\item \textbf{The causes are new.} The polluter search ran on 62 of the 65 flaky held out Claude tests (kept or removed). It (\texttt{od\_diagnose.py}) found a single polluter for 20, 10 failed alone (7 only when the hash seed varies), and 32 passed alone and after every single other test. 20 of those 32 come from one Sonnet kindle-maker suite. Its fixture makes a folder named \texttt{/tmp/kmtest\_<id(self)>\_1}, where \texttt{id()} is a memory address, and a later test can get the same address while a folder from a test that did not clean up still exists. These tests depend on the order through the file system, which a pairwise search inside one process cannot reproduce. In fistro, Haiku and Sonnet tests call an unseeded date generator that can pick a day the month does not have, and in jmcm a Haiku suite deletes the package from \texttt{sys.modules} and imports it again, the habit that made GPT~5.4 mini flaky in Progress Report 2. These are machine findings from the reruns. Manas will label every new flaky test with the method of Section~\ref{sec:rootcause3}.
\end{itemize}

\textbf{The pre registered test.} \PENDING{PREREG}

\subsection{RQ1, prevalence: order dependent and time batches}
Table~\ref{tab:batches} gives the reproduction results of the two batches of Progress Report 2's steps 3 and 4.

\begin{table}[h]
\caption{Reproducing the two new batches of Gruber et al.\ flaky tests}
\label{tab:batches}
\small
\begin{tabular}{lrr}
\toprule
Outcome & Time & Order dependent \\
\midrule
Tried & 16 & 45 \\
\quad Repository gone or install failed & 0 & 5 \\
\quad Always failed & 12 & 17 \\
\quad Never failed & 4 & 6 \\
\quad Reproduced & 0 & 17 \\
\bottomrule
\end{tabular}
\end{table}

\textbf{Time batch.} Only 16 tests in the unlabeled data passed the rule, and none reproduced: 12 always failed, mostly at collection on old dependencies, and 4 never failed in 150 runs. The known flaky code study therefore still has only six projects outside randomness, which is a limit we state rather than fix (Section~\ref{sec:threats3}).

\textbf{Order dependent arm.} 17 of 45 order dependent tests reproduced (38\%, well above the 22\% yield of the earlier batches), so we generated 4 suites per Claude model for each of the 17 projects (136 suites) and ran them like the main study. Run alone, 11 of the 17 developer tests always passed and 5 always failed, so most are victims and 5 are brittles that need another test to set state first. Only one test, ChainConsumer, flipped alone (13 of 200). The generated tests behave differently. Haiku wrote 1,644 kept tests with none flaky, and Sonnet 1,403 with 3 kept flaky tests, all in HikLoad and all shuffle only. They reimport \texttt{hikload.config} and check that the module's \texttt{CONFIG} dict is the same object, which fails when another test has reloaded the module. So on code where the developer tests depend on order, the generated tests almost never do (3 of 3,047 kept tests). \PENDING{GPTOD}

\subsection{RQ3, features of flaky tests}
\label{sec:rq3}
Table~\ref{tab:features} extends the Progress Report 2 feature table to randomness, time, and concurrency, and adds for each feature how many of the tests that use it flipped. It covers the library study, where human and LLM tests ran on the same libraries (\texttt{scripts/feature\_flakiness.py}).

\begin{table}[h]
\caption{Library study test functions with each feature: how many use it (share of all functions in \%) and how many of those flipped. Human: 3,077 functions, 4 flaky. LLM: 7,474 functions in all seven conditions, 8 flaky.}
\label{tab:features}
\small
\setlength{\tabcolsep}{2.5pt}
\begin{tabularx}{\columnwidth}{Xrrrr}
\toprule
 & \multicolumn{2}{c}{Human} & \multicolumn{2}{c}{LLM} \\
Feature & Use (\%) & Flaky & Use (\%) & Flaky \\
\midrule
Calls a random generator & 46 (1.5) & 0 & 138 (1.8) & 1 \\
\quad without a seed & 34 (1.1) & 0 & 82 (1.1) & 1 \\
Reads the real clock & 58 (1.9) & 0 & 206 (2.8) & 3 \\
Sleeps for real & 35 (1.1) & 1 & 107 (1.4) & 0 \\
Asserts on elapsed time & 4 (0.1) & 0 & 42 (0.6) & 1 \\
Threads, processes, or asyncio & 56 (1.8) & 1 & 248 (3.3) & 0 \\
\quad thread started without join & 2 (0.1) & 0 & 6 (0.1) & 0 \\
Shares state between tests & 736 (23.9) & 2 & 170 (2.3) & 0 \\
Class or module scoped fixture & 599 (19.5) & 2 & 0 & 0 \\
Setup without teardown & 130 (4.2) & 1 & 0 & 0 \\
Freezes the clock (freezegun) & 12 (0.4) & 0 & 323 (4.3) & 0 \\
Uses \texttt{monkeypatch} & 7 (0.2) & 0 & 133 (1.8) & 3 \\
\bottomrule
\end{tabularx}
\end{table}

\textbf{Humans use the risky features and stay stable.} 34 human test functions use randomness without a seed and none flipped. 58 read the real clock and none flipped, and 55 of 56 that use threads, processes, or asyncio never flipped. The same holds for LLM tests: 1 of 82 unseeded randomness functions and 3 of 206 clock reading functions flipped. A risky feature is a weak predictor in both groups, which matches the scanner precision of 22\% in Progress Report 2. \textbf{The groups use the features differently.} LLM tests read the clock, assert on elapsed time, and start threads more often than human tests, but they also freeze the clock ten times as often (4.3\% against 0.4\%) and use \texttt{monkeypatch} nine times as often. Human tests share state through wide fixtures, and both human order dependent flips (arrow) sit on such a fixture. No LLM test uses one. The 3 LLM tests with \texttt{monkeypatch} that flipped are the humanize brittles of Progress Report 2, where the patch itself is the problem. With 4 and 8 flaky functions the per feature rates have wide intervals, so the table describes use, not risk.

\subsection{RQ2, root causes}
\label{sec:rootcause3}
The 91 labels of Progress Report 2 are unchanged. For the agreement check, \texttt{results/root\_causes/blind/sheet.csv} lists 78 of the labeled tests in a shuffled order with the model and Manas's label hidden. Shaunak fills it in without seeing those, and \texttt{scripts/kappa.py} then computes Cohen's kappa. The new flaky tests of the held out sample and the order dependent arm will be labeled the same way, with the polluter search results above as evidence.

\subsection{RQ4, mitigation}
\label{sec:rq4}
\textbf{Fixtures on tests they were not built for.} M1 and M2 were frozen at commit \texttt{7c05b6b} and ran unchanged on every held out suite (Table~\ref{tab:fixtures}). This is the prospective check Progress Report 2 asked for.

\begin{table}[h]
\caption{Kept flaky tests in the held out sample and the order dependent arm, plain and with each fixture, 200 runs each}
\label{tab:fixtures}
\small
\setlength{\tabcolsep}{3pt}
\begin{tabularx}{\columnwidth}{Xrrr}
\toprule
Suites & Plain & M1, seed & M2, seed and restore \\
\midrule
Sonnet, held out & 18 & 14 & 12 \\
Haiku, held out & 34 & 20 & 3 \\
%GPTFIX%
Sonnet, order dependent arm & 3 & 3 & 0 \\
Haiku, order dependent arm & 0 & 0 & 0 \\
\bottomrule
\end{tabularx}
\end{table}

\begin{itemize}
\item \textbf{M2 generalizes to new order dependent tests in memory.} Haiku goes from 34 kept flaky tests to 3, including the 17 jmcm tests that delete the package from \texttt{sys.modules}, and the 3 HikLoad tests of the order dependent arm stop flipping. M1 alone removes only the randomness cases.
\item \textbf{It does not reach state outside the interpreter.} Sonnet keeps 12 under M2, and 11 are the kindle-maker tests that collide through \texttt{/tmp}. M2 restores modules, globals, the environment, and the working directory, not the file system.
\item \textbf{No stable test broke.} M2 made 3 stable tests fail in Progress Report 2. On the 13,879 kept held out Claude tests and the 3,047 of the order dependent arm, no test that passed every plain run failed in every M2 run. The keplertools failure is still unexplained (Section~\ref{sec:challenges3}).
\end{itemize}

\textbf{Guard v3 on two new models.} We generated 2 guard v3 suites per Claude model for all 41 projects and compare them with held out samples 1 and 2, which come from the same pipeline (Table~\ref{tab:guard3}). \PENDING{GPTGUARD}

\begin{table}[h]
\caption{Guard v3 against the plain prompt, same pipeline, samples 1 and 2 of each, 41 projects. Coverage is the mean line coverage of the target files, mutation the pooled mutation score.}
\label{tab:guard3}
\small
\setlength{\tabcolsep}{2.5pt}
\begin{tabularx}{\columnwidth}{Xrrrrr}
\toprule
Suites & No import & Kept & Flaky kept & Cov. & Mut. \\
\midrule
Sonnet plain & 0 & 1,735 & 2 & 79.0 & 60.1 \\
Sonnet guard v3 & 0 & 1,622 & 0 & 78.4 & 61.1 \\
Haiku plain & 1 & 1,840 & 21 & 74.9 & 52.8 \\
Haiku guard v3 & 2 & 1,616 & 0 & 73.7 & 53.2 \\
%GPTGUARDROWS%
\bottomrule
\end{tabularx}
\end{table}

For the Claude models guard v3 removed every kept flaky test (Haiku 21 in 5 suites to 0 in 80 suites, Fisher's exact test on suites $p = 0.059$) and cost 6 to 12\% of the kept tests, with no loss of coverage or mutation score. That is unlike GPT~4.1 mini, where the guard prompts made suites fail to import (Table~\ref{tab:quality}).

\textbf{Test quality.} Table~\ref{tab:quality} answers the question whether the mitigations cost quality. Every number uses the same mutants per project for every suite.

\begin{table}[h]
\caption{Quality of the GPT~4.1 mini guard prompts on the first 17 projects (34 suites each). Mut., imported is the mutation score over the suite pairs where both suites import.}
\label{tab:quality}
\small
\setlength{\tabcolsep}{2.5pt}
\begin{tabularx}{\columnwidth}{Xrrrrr}
\toprule
Prompt & No kept & Kept & Cov. & Mut. & Mut., imported \\
\midrule
Plain & 1 & 559 & 80.9 & 55.7 & \\
Guard v1 & 1 & 548 & 81.5 & 49.8 & 58.1 to 51.9 \\
Guard v2 & 9 & 500 & 57.7 & 36.2 & 56.6 to 55.3 \\
Guard v3 & 7 & 428 & 66.9 & 40.5 & 55.2 to 50.4 \\
\bottomrule
\end{tabularx}
\end{table}

For GPT~4.1 mini most of the cost comes from suites that no longer import. When both suites import, coverage is unchanged and the mutation score drops 1 to 6 points. Across the 41 projects with the plain prompt, the Claude suites are also stronger than the GPT suites, with mean coverage 79 against 73\% and mutation scores of 61 to 62 against 40 to 41\%, so a model that writes fewer flaky tests is not buying that with weaker tests. Fixtures do not change the tests, so their only cost is the tests they break, which was none in the held out sample.

\textbf{Measured rerun filters.} The runner now records every outcome in run order, so we reran the 14 projects with flaky suites 200 more times and measured the filters directly instead of estimating them (\texttt{scripts/krun\_measured.py}). 74 tests flipped. A filter that reruns the file in its written order keeps 57 of them after 1 run, 49 after 3, 42 after 5, and 38 after 10. 32 of the 74 never fail in file order, so no number of file order reruns catches them. Of the 42 that do fail in file order, 25 pass the first run and 6 still pass 10. A filter that alternates file order and shuffled runs keeps 28 after 10 runs. This confirms the estimate of Progress Report 2 with measured data: extra reruns help only if they change the order.

\subsection{Threats to validity}
\label{sec:threats3}
\begin{itemize}
\item \textbf{Scope.} 35 of 41 main projects are about randomness, the time batch reproduced nothing, and the order dependent arm has 17 projects. We cannot claim anything about time or concurrency flakiness in LLM tests beyond the library study.
\item \textbf{Pipeline.} The held out sample removes the pipeline difference of Progress Report 2 for the comparison it tests, but the earlier suites still mix two pipelines, so we do not pool them with it.
\item \textbf{Clustering.} 11 of the 52 kept flaky held out Claude tests come from one kindle-maker suite and 17 from one jmcm suite. We count suites and projects for every comparison.
\item \textbf{Our own mitigations.} M2 and guard v3 were tuned on the first 17 projects, so we report their held out results apart from the tuning results.
\end{itemize}


\section{Milestone goals and status}
Table~\ref{tab:goals} lists the five steps Progress Report 2 promised for this milestone, plus the check point rules it set.

\begin{table*}[t]
\caption{Goals for this milestone, from the next steps of Progress Report 2}
\label{tab:goals}
\small
\begin{tabularx}{\textwidth}{p{5.6cm}p{1.6cm}X}
\toprule
Goal & Status & Notes \\
\midrule
1. Held out sample: 8 new plain suites per model for the 41 projects in one pipeline, run plain, with M1, and with M2, analyzed with the frozen sign test & \PENDING{G1S} & \PENDING{G1N} \\
2. Mitigation for every model: guard v3 for the other three models, a measured 3 and 5 run filter, why M2 broke 3 tests, M1 and M2 on the library study & Mostly done & Guard v3 ran for Sonnet and Haiku\PENDING{G2N}. The rerun filters are now measured from outcome sequences. M2 broke no held out test, but the keplertools cause is still unknown. M1 and M2 on the library study and the comparison with an iFixFlakies style fix moved to the final report. \\
3. Order dependent arm, 10 projects (5 if short) & Done & 17 of 45 reproduced. 4 suites per model per project\PENDING{G3N}. We ran the generated suites on their own. Running them inside the developer suite, to see if they pollute developer tests, is a next step. \\
4. 40 more time, concurrency, and async wait tests & Done, no yield & Only 16 passed the rule and none reproduced. By the check point rule we did not run stress-ng on the library study instead, because the order dependent arm took the time. \\
5. Root causes and a blind relabel with Cohen's kappa & In progress & The blind sheet (78 tests) and \texttt{kappa.py} are ready for Shaunak's relabel. The new flaky tests have polluter search results and go to Manas for labels. \\
New: test quality after mitigation (Progress Report 2 feedback) & Done & Coverage and mutation score for every plain and guard suite (Tables~\ref{tab:guard3} and~\ref{tab:quality}). \\
\bottomrule
\end{tabularx}
\end{table*}

The final report target was at least 40 suites with a kept flaky test from at least three models, and at least 8 projects outside randomness. We now have \PENDING{SUITECOUNT}, so the first target is met. The second is not, because the time batch reproduced nothing, and we report that as a limit.

\section{Challenges}
\label{sec:challenges3}
\begin{itemize}
\item \textbf{Reproduction yield.} 0 of 16 time tests reproduced, while 17 of 45 order dependent tests did. Old time tests mostly fail at collection on dependencies that no longer install.
\item \textbf{Credentials.} The Claude CLI login and the OpenAI key both expired during the milestone, which delayed the GPT half of the held out sample by about two days. All generation now runs through a job runner that skips finished suites and stops at a spending cap (\texttt{scripts/gen\_jobs.py}).
\item \textbf{Slow projects.} The PyTorch projects take close to the four hour limit of a GitHub job for 16 suites of 200 runs.
\item \textbf{The M2 keplertools failure.} Two stable tests compute an eccentricity of 0 through \texttt{np.sqrt(1.0 - p/a)}, which gives NaN when rounding makes the argument slightly negative. Under M2 this happens on the Linux runners and not on our laptop. The project has no fast math build flags, so we have not found what M2 changes.
\item \textbf{File system state.} The kindle-maker tests show that a test can depend on order through files a previous test left behind, which neither our pairwise polluter search nor M2 sees.
\end{itemize}

\section{Next steps toward the final report}
\label{sec:next3}
The final report is due December 12. In priority order:
\begin{enumerate}
\item \textbf{Finish the held out analysis (Shaunak, by November 22).} \PENDING{NEXT1}
\item \textbf{Root causes and agreement (Manas, by November 25; Shaunak's relabel by November 22).} Label every new flaky test, and report Cohen's kappa on the blind sheet.
\item \textbf{Order dependent arm inside the developer suite (Manas and Shaunak, by November 29).} Run each generated suite inside the project's developer suite in shuffled orders, to test whether generated tests become polluters of developer tests, which Progress Report 2 promised.
\item \textbf{Fixtures on the library study and a file system restore (Shaunak, by December 2).} Run M1 and M2 on the library study suites, add a temporary directory reset to M2 for the kindle-maker case, and compare with an iFixFlakies style cleaner.
\item \textbf{Writing (both, by December 10).} Manas writes related work and the root cause section, Shaunak the setup, results, and threats. Section~\ref{sec:setup} is our first draft of the setup.
\end{enumerate}

\textbf{Division of work.} Since Progress Report 2, Shaunak built the held out pipeline, generated and ran the held out, guard v3, and order dependent arm suites, wrote the quality, feature, and rerun filter analyses, and wrote this report. \PENDING{MANAS}

\section{Code, data, and documentation}
\label{sec:code}
Everything is in the public repository \url{https://github.com/shaunak1206/ece382v-flaky-llm-tests}, with this report's results at commit \texttt{\PENDING{COMMIT}}. \texttt{docs/DATA.md} describes every result file. The held out analysis is \texttt{python scripts/heldout\_stats.py} (the frozen sign test) and \texttt{python scripts/heldout\_summary.py}, the quality tables \texttt{python scripts/quality\_summary.py}, the rerun filters \texttt{python scripts/krun\_measured.py}, and the feature table \texttt{python scripts/feature\_flakiness.py}, all offline from the committed results.

\bibliographystyle{ACM-Reference-Format}
\bibliography{refs}

\end{document}
